\chapter{Минимальный смешивающий коннектор на фиксированном \texorpdfstring{$H_{15}$}{H15}}
\label{ch:v7-minimal-h15-mixed-connector-admission-gate}

\section{Вопрос и строгий вход}

Предыдущие гейты доказали, что физическая вторая степень и все положительные
односледовые моменты текущего оператора распадаются по кромкам $u,d,e$.
Поэтому новый потенциал от того же блочного дерева не допускается. Требуется
классифицировать минимальную новую стрелочную структуру, впервые замыкающую
смешанный цикл.

Фиксируем фермионные вершины заряженного $H_{15}$:
\begin{equation}
 L=\{Q_L,L_L\},\qquad R=\{u_R,d_R,e_R\},
\end{equation}
и существующие рёбра
\begin{equation}
 E_0=\{Q_Lu_R,Q_Ld_R,L_Le_R\}.
 \label{eq:v7-existing-h15-forest}
\end{equation}
Это лес без четырёхцикла. Старый гейт Тома IV установил, что в
бимодульном квадрате первым order-one-совместимым циклом может быть
четырёхкромочный прямоугольник. Здесь это используется только как
предклассификация: полный Real/order-one тест конкретного расширения ещё
не выполнен.

\section{Полный перебор прямоугольных дополнений}

На фиксированных вершинах отсутствуют ровно три рёбра:
\begin{equation}
 E_{\rm miss}=\{L_Lu_R,L_Ld_R,Q_Le_R\}.
\end{equation}
Одной новой стрелки недостаточно. Минимум равен двум, и полный перебор даёт
ровно три дополнения:
\begin{align}
 \{L_Lu_R,L_Ld_R\}&:\quad \{u_R,d_R\},\nonumber\\
 \{L_Lu_R,Q_Le_R\}&:\quad \{u_R,e_R\},
 \label{eq:v7-three-minimal-rectangles}\\
 \{L_Ld_R,Q_Le_R\}&:\quad \{d_R,e_R\}.\nonumber
\end{align}
Правая пара обозначает два столбца соответствующего $K_{2,2}$.

\section{Калибровочная типизация}

Используем стандартные представления
\begin{align}
 Q_L&\sim(\mathbf3,\mathbf2)_{1/6},&
 L_L&\sim(\mathbf1,\mathbf2)_{-1/2},\nonumber\\
 u_R&\sim(\mathbf3,\mathbf1)_{2/3},&
 d_R&\sim(\mathbf3,\mathbf1)_{-1/3},&
 e_R&\sim(\mathbf1,\mathbf1)_{-1}.
\end{align}
Для вершины $\overline\psi_L S\psi_R$ калибровочная инвариантность требует
\begin{equation}
 Y_S=Y_L-Y_R.
 \label{eq:v7-scalar-hypercharge-rule}
\end{equation}
Поэтому недостающие рёбра имеют типы
\begin{align}
 L_Lu_R&:\quad S\sim(\overline{\mathbf3},\mathbf2)_{-7/6},\nonumber\\
 L_Ld_R&:\quad S\sim(\overline{\mathbf3},\mathbf2)_{-1/6},
 \label{eq:v7-missing-edge-scalars}\\
 Q_Le_R&:\quad S\sim(\mathbf3,\mathbf2)_{7/6}.\nonumber
\end{align}
Первое и третье являются сопряжёнными чтениями одного комплексного
мультиплета
\begin{equation}
 R_2\sim(\mathbf3,\mathbf2)_{7/6}.
 \label{eq:v7-r2-candidate}
\end{equation}
Он поддерживает два блока
\begin{equation}
 \overline Q_LR_2e_R,
 \qquad
 \overline u_RR_2^Ti\sigma_2L_L+\mathrm{h.c.}
 \label{eq:v7-r2-two-edges}
\end{equation}
и единственный замыкает минимальный прямоугольник $\{u_R,e_R\}$ одним
комплексным полем. Два других дополнения требуют одновременно $R_2$ и
независимый $\widetilde R_2\sim(\mathbf3,\mathbf2)_{1/6}$.

\section{Почему второй обычный Хиггс не решает этот гейт}

Второй дублет $(\mathbf1,\mathbf2)_{1/2}$ может дать независимые поля на
уже существующих рёбрах, но не меняет $E_0$, не создаёт отсутствующую
матричную единицу между метками $u,d,e$ и не замыкает $K_{2,2}$. Это
отдельная расширенная скалярная ветвь, а не требуемый смешанный коннектор.

\section{Литературная и физическая граница}

Диаграммная классификация конечных спектральных троек связывает вершины с
хиральными калибровочными мультиплетами, а стрелки --- с юкавскими блоками
\cite{v7connector-krajewski}. Лептокварковые блоки в ранней NCG-модели
появлялись после отказа от $S^0$-реальности, причём конкретная реализация
не прошла феноменологические ограничения \cite{v7connector-paschke}.
Clifford-расширение также порождает скаляры одновременно со слабым и
цветовым индексами \cite{v7connector-kurkov}. Двуххиггсовская NCG-модель
существует как отдельная архитектура, но сама по себе не создаёт нужный
кварк-лептонный блок \cite{v7connector-jimenez}.

Следовательно, \eqref{eq:v7-r2-candidate} --- единственный графово
минимальный кандидат на прежнем фермионном наборе, но не элемент текущего
одноформенного бимодуля $\mathbb C^3$. До физического допуска необходимы:
\begin{enumerate}
 \item полный тест первого порядка и Real/$J$-совместимости;
 \item вывод обоих блоков \eqref{eq:v7-r2-two-edges} одной флуктуацией;
 \item цветосохраняющий вакуум полного спектрального действия;
 \item аудит недопустимых фермионных и барионных каналов;
 \item семейно-чувствительное смешанное слово и полный гессиан.
\end{enumerate}

\begin{theorem}[Минимальность смешанного кандидата]
На фиксированных фермионных вершинах заряженного $H_{15}$ первый
четырёхкромочный смешанный цикл требует не менее двух новых рёбер. Существует
ровно три минимальных дополнения. Единственное дополнение одним комплексным
калибровочным мультиплетом замыкает пару $u_R,e_R$ посредством
$R_2\sim(\mathbf3,\mathbf2)_{7/6}$ и его сопряжения.
\end{theorem}

\begin{nogocriterion}[Классификация не равна допуску]
$R_2$ не принадлежит неизменённому родителю и не считается выведенным
полем. Пока не пройдены Real, первое условие, спектральное действие и
цветовой вакуум, это только классификация минимального расширения.
\end{nogocriterion}

\begin{protocol}
\begin{enumerate}
 \item четырёхцикл в текущем $H_{15}$: \textbf{нет};
 \item минимальное число новых рёбер: \textbf{два};
 \item минимальных дополнений: \textbf{три};
 \item одно-мультиплетное дополнение: \textbf{единственно};
 \item его тип: \textbf{$R_2=(\mathbf3,\mathbf2)_{7/6}$};
 \item второй обычный Хиггс замыкает граф: \textbf{нет};
 \item $R_2$ содержится в текущем модуле: \textbf{нет};
 \item физический допуск: \textbf{условно открыт}.
\end{enumerate}
\end{protocol}

Машинный сертификат сохранён в
\path{s2t_v7_minimal_h15_mixed_connector_admission_gate_results.json}.

\begin{thebibliography}{9}
\bibitem{v7connector-krajewski}
T.~Krajewski, \emph{Classification of Finite Spectral Triples},
Journal of Geometry and Physics 28 (1998), 1--30; arXiv:hep-th/9701081.
\bibitem{v7connector-paschke}
M.~Paschke, F.~Scheck, A.~Sitarz,
\emph{Can (noncommutative) geometry accommodate leptoquarks?},
Physical Review D 59 (1999), 035003; arXiv:hep-th/9709009.
\bibitem{v7connector-kurkov}
M.~A.~Kurkov, F.~Lizzi,
\emph{Clifford Structures in Noncommutative Geometry and the Extended Scalar Sector},
Physical Review D 97 (2018), 085024; arXiv:1801.00260.
\bibitem{v7connector-jimenez}
F.~Jimenez, D.~Restrepo, A.~Rivera,
\emph{Type-II two-Higgs-doublet model in noncommutative geometry},
arXiv:2202.04222.
\end{thebibliography}